\paragraph{双指纹审计的张量闭合、最小二值统计与强乘法律}
在 \(P_{10}\times L_{\mathrm{LY}}\) 的双域 Chebotarev 乘积律、三域类函数张量分解、\(\chi_{\mathrm{LY}}\) 的唯一有效位以及 Latt\`es 六次域的统一稀释律已经确立之后，还可把双指纹审计进一步压缩为下列五条闭合后果。它们依次对应标准素数平均下的类函数张量因子化、一次因子事件与 Lee--Yang 三次的三种分裂型的完整联合密度、由“是否有根”与判别式符号构成的最小二值统计量、仅保留单一符号位时的精确信息亏损，以及任意有限布尔组合上的强乘法律。

\begin{theorem}[标准素数平均下的双指纹类函数张量因子化]\label{thm:xi-time-part62ea-dualfingerprint-classfunction-prime-average-factorization}
\leanverified{paper\_xi\_time\_part62ea\_dualfingerprint\_classfunction\_prime\_average\_factorization}
沿用命题 \ref{prop:fold-gauge-anomaly-P10-leyang-linear-disjointness} 的记号。对任意类函数
\[
f:S_{10}\to\CC,
\qquad
h:S_3\to\CC,
\]
对每个共同不分歧素数 \(p\) 记
\[
\sigma_{10}(p):=\mathrm{Frob}_p(L_{10}/\QQ)\in S_{10},
\qquad
\sigma_3(p):=\mathrm{Frob}_p(L_{\mathrm{LY}}/\QQ)\in S_3.
\]
则有极限恒等式
\[
\lim_{x\to\infty}
\frac{1}{\pi(x)}
\sum_{\substack{p\le x\\ p\nmid \Disc(P_{10})\Disc(P_{\mathrm{LY}})}}
f\!\bigl(\sigma_{10}(p)\bigr)\,
h\!\bigl(\sigma_3(p)\bigr)
=
\left(\frac{1}{10!}\sum_{\tau\in S_{10}}f(\tau)\right)
\left(\frac{1}{6}\sum_{\upsilon\in S_3}h(\upsilon)\right).
\]
特别地，若定义中心化类函数
\[
\widetilde f:=
f-\frac{1}{10!}\sum_{\tau\in S_{10}}f(\tau),
\qquad
\widetilde h:=
h-\frac{1}{6}\sum_{\upsilon\in S_3}h(\upsilon),
\]
则
\[
\lim_{x\to\infty}
\frac{1}{\pi(x)}
\sum_{\substack{p\le x\\ p\nmid \Disc(P_{10})\Disc(P_{\mathrm{LY}})}}
\widetilde f\!\bigl(\sigma_{10}(p)\bigr)\,
\widetilde h\!\bigl(\sigma_3(p)\bigr)
=0.
\]
\end{theorem}
\begin{proof}
由定理 \ref{thm:fold-gauge-anomaly-P10-leyang-class-function-tensor-independence}，
\[
\lim_{x\to\infty}
\frac{1}{\pi_{10,\mathrm{LY}}^\ast(x)}
\sum_{\substack{p\le x\\ p\nmid \Disc(P_{10})\Disc(P_{\mathrm{LY}})}}
f\!\bigl(\sigma_{10}(p)\bigr)\,
h\!\bigl(\sigma_3(p)\bigr)
=
\left(\frac{1}{10!}\sum_{\tau\in S_{10}}f(\tau)\right)
\left(\frac{1}{6}\sum_{\upsilon\in S_3}h(\upsilon)\right),
\]
其中
\[
\pi_{10,\mathrm{LY}}^\ast(x):=
\#\{p\le x:\ p\nmid \Disc(P_{10})\Disc(P_{\mathrm{LY}})\}.
\]
由于坏素数只有有限多个，
\[
\frac{\pi_{10,\mathrm{LY}}^\ast(x)}{\pi(x)}\longrightarrow 1
\qquad(x\to\infty),
\]
故把分母由 \(\pi_{10,\mathrm{LY}}^\ast(x)\) 改为 \(\pi(x)\) 不改变极限值，从而得到第一式。

第二式只需把 \(f,h\) 分别替换为 \(\widetilde f,\widetilde h\)。此时两者的群平均均为 \(0\)，故右端为 \(0\)。证毕。
\end{proof}

\begin{corollary}[一次因子事件与 Lee--Yang 三种分裂型的完整联合密度]\label{cor:xi-time-part62ea-A-leyang-threeway-density}
沿用定理 \ref{thm:xi-time-part62ea-dualfingerprint-classfunction-prime-average-factorization} 的记号，并令
\[
g(y):=P_{\mathrm{LY}}(y)=256y^3+411y^2+165y+32.
\]
对共同不分歧素数 \(p\)，定义事件
\[
A:=\{\text{\(P_{10}\bmod p\) 含一次因子}\},
\]
\[
B_{\mathrm{split}}:=\{\text{\(g\bmod p\) 全分裂}\},
\qquad
B_{\mathrm{root}}:=\{\text{\(g\bmod p\) 至少有一根}\},
\]
\[
B_{\mathrm{single}}:=B_{\mathrm{root}}\setminus B_{\mathrm{split}},
\qquad
B_{\mathrm{noroot}}:=B_{\mathrm{root}}^{c}.
\]
则有
\[
\delta(A\cap B_{\mathrm{split}})=\frac{28319}{268800},
\]
\[
\delta(A\cap B_{\mathrm{single}})=\frac{28319}{89600},
\]
\[
\delta(A\cap B_{\mathrm{noroot}})=\frac{28319}{134400}.
\]
特别地，
\[
\delta(A\mid B_{\mathrm{split}})
=
\delta(A\mid B_{\mathrm{single}})
=
\delta(A\mid B_{\mathrm{noroot}})
=
\frac{28319}{44800}.
\]
\end{corollary}
\begin{proof}
由推论 \ref{cor:fold-gauge-anomaly-P10-modp-splitting-densities}，
\[
\delta(A)=\frac{28319}{44800}.
\]
由结论 \ref{con:xi-terminal-zm-leyang-cubic-modp-splitting-density}，
\[
\delta(B_{\mathrm{split}})=\frac16,
\qquad
\delta(B_{\mathrm{single}})=\frac12,
\qquad
\delta(B_{\mathrm{noroot}})=\frac13.
\]
再由推论 \ref{cor:fold-gauge-anomaly-P10-leyang-chebotarev-product-law}，\(A\) 与每个只依赖 \(S_3\) 分量的事件都严格独立，故
\[
\delta(A\cap E)=\delta(A)\delta(E)
\]
对 \(E\in\{B_{\mathrm{split}},B_{\mathrm{single}},B_{\mathrm{noroot}}\}\) 成立。代入上述边际密度即得三条联合密度公式。最后对各式分别除以 \(\delta(E)\) 便得到条件密度恒等式。证毕。
\end{proof}

\begin{theorem}[Lee--Yang 三分裂型的最小二值充分统计量]\label{thm:xi-time-part62ea-leyang-minimal-binary-sufficient-statistic}
\leanverified{paper\_xi\_time\_part62ea\_leyang\_minimal\_binary\_sufficient\_statistic}
沿用推论 \ref{cor:xi-time-part62ea-A-leyang-threeway-density} 的记号，并记
\[
\Sigma_{\mathrm{LY}}(p)\in\{(1)(1)(1),(1)(2),(3)\}
\]
为 \(g\bmod p\) 的分解型随机变量。定义二值观测对
\[
\mathcal O(p):=
\bigl(\mathbf 1_{B_{\mathrm{root}}}(p),\chi_{\mathrm{LY}}(p)\bigr)
\in\{0,1\}\times\{\pm1\}.
\]
则 \(\mathcal O\) 的像恰为
\[
\{(1,+1),(1,-1),(0,+1)\},
\]
并且存在双射
\[
(1)(1)(1)\longleftrightarrow (1,+1),
\qquad
(1)(2)\longleftrightarrow (1,-1),
\qquad
(3)\longleftrightarrow (0,+1).
\]
因而 \((0,-1)\) 从不出现，且 \(\mathcal O\) 与 \(\Sigma_{\mathrm{LY}}\) 等价地编码同一三值信息。特别地，在一切由二值坐标组成并能逐素完全恢复 \(\Sigma_{\mathrm{LY}}\) 的观测中，\(\mathcal O\) 具有最小坐标数；因此它构成 Lee--Yang 三分裂型的最小充分二值统计量。
\end{theorem}
\begin{proof}
由结论 \ref{con:xi-terminal-zm-leyang-cubic-modp-splitting-density}，
\[
\Sigma_{\mathrm{LY}}=(1)(1)(1)
\iff
g\bmod p\ \text{全分裂},
\]
\[
\Sigma_{\mathrm{LY}}=(1)(2)
\iff
g\bmod p\ \text{恰有一个一次因子},
\]
\[
\Sigma_{\mathrm{LY}}=(3)
\iff
g\bmod p\ \text{在 }\FF_p\text{ 上无根}.
\]
另一方面，由定理 \ref{thm:xi-time-part62e-chi-ly-sole-controller}，
\[
\chi_{\mathrm{LY}}(p)=-1
\iff
\Sigma_{\mathrm{LY}}(p)=(1)(2),
\]
\[
\chi_{\mathrm{LY}}(p)=+1
\iff
\Sigma_{\mathrm{LY}}(p)\in\{(1)(1)(1),(3)\}.
\]
再由 \(\mathbf 1_{B_{\mathrm{root}}}(p)=1\) 当且仅当 \(\Sigma_{\mathrm{LY}}(p)\in\{(1)(1)(1),(1)(2)\}\)，立刻得到上述对应表，故 \(\mathcal O\) 与 \(\Sigma_{\mathrm{LY}}\) 双射，因而充分性成立。

若只允许一个二值坐标，则观测值至多只有 \(2\) 种，而 \(\Sigma_{\mathrm{LY}}\) 有 \(3\) 种可能状态，因此不可能由单一二值观测完全恢复。故任何完备恢复 \(\Sigma_{\mathrm{LY}}\) 的二值编码都至少需要两个坐标；而 \(\mathcal O\) 恰以两个二值坐标实现该恢复，故其为最小。证毕。
\end{proof}

\begin{corollary}[单一符号位压缩的精确信息亏损]\label{cor:xi-time-part62ea-chi-ly-exact-information-loss}
沿用定理 \ref{thm:xi-time-part62ea-leyang-minimal-binary-sufficient-statistic} 的记号，并以 \(\log_2\) 为底定义 Shannon 熵。则有
\[
H(\Sigma_{\mathrm{LY}}\mid \chi_{\mathrm{LY}})
=
\frac12\,h_2\!\left(\frac13\right),
\qquad
h_2(x):=-x\log_2 x-(1-x)\log_2(1-x).
\]
等价地，
\[
H(\Sigma_{\mathrm{LY}})
=
1+\frac12\,h_2\!\left(\frac13\right),
\qquad
I(\Sigma_{\mathrm{LY}};\chi_{\mathrm{LY}})=1.
\]
数值上，
\[
H(\Sigma_{\mathrm{LY}}\mid \chi_{\mathrm{LY}})
\approx 0.4591479170\ \text{bits}.
\]
\end{corollary}
\begin{proof}
由定理 \ref{thm:killo-leyang-quadratic-character-decoupling}，
\[
\delta(\chi_{\mathrm{LY}}=+1)=\delta(\chi_{\mathrm{LY}}=-1)=\frac12.
\]
再由定理 \ref{thm:xi-time-part62ea-leyang-minimal-binary-sufficient-statistic}，
\[
\chi_{\mathrm{LY}}=-1
\Longrightarrow
\Sigma_{\mathrm{LY}}=(1)(2),
\]
故
\[
H(\Sigma_{\mathrm{LY}}\mid \chi_{\mathrm{LY}}=-1)=0.
\]
而在 \(\chi_{\mathrm{LY}}=+1\) 的条件下，只剩两种可能：
\[
\Sigma_{\mathrm{LY}}=(1)(1)(1)\quad\text{或}\quad \Sigma_{\mathrm{LY}}=(3).
\]
其条件概率由 \(S_3\) 中恒等类与三循环类的类大小比 \(1:2\) 给出，即
\[
\mathbb P\!\bigl(\Sigma_{\mathrm{LY}}=(1)(1)(1)\mid \chi_{\mathrm{LY}}=+1\bigr)=\frac13,
\]
\[
\mathbb P\!\bigl(\Sigma_{\mathrm{LY}}=(3)\mid \chi_{\mathrm{LY}}=+1\bigr)=\frac23.
\]
因此
\[
H(\Sigma_{\mathrm{LY}}\mid \chi_{\mathrm{LY}}=+1)=h_2\!\left(\frac13\right).
\]
对 \(\chi_{\mathrm{LY}}=\pm1\) 按各自概率 \(1/2\) 加权，得到
\[
H(\Sigma_{\mathrm{LY}}\mid \chi_{\mathrm{LY}})
=
\frac12\,h_2\!\left(\frac13\right).
\]

又由于 \(\chi_{\mathrm{LY}}\) 是 \(\Sigma_{\mathrm{LY}}\) 的确定函数且取值等分布，故
\[
H(\chi_{\mathrm{LY}})=1,
\qquad
I(\Sigma_{\mathrm{LY}};\chi_{\mathrm{LY}})=H(\chi_{\mathrm{LY}})=1.
\]
将链式公式
\[
H(\Sigma_{\mathrm{LY}})
=
I(\Sigma_{\mathrm{LY}};\chi_{\mathrm{LY}})
+H(\Sigma_{\mathrm{LY}}\mid \chi_{\mathrm{LY}})
\]
代入即可得到第二式。证毕。
\end{proof}

\begin{theorem}[双指纹有限布尔组合的强乘法律]\label{thm:xi-time-part62ea-dualfingerprint-boolean-strong-product-law}
\leanverified{paper\_xi\_time\_part62ea\_dualfingerprint\_boolean\_strong\_product\_law}
仍在共同不分歧素数的密度概率空间上，令 \(\mathcal A_{10}\) 为由全部事件
\[
\{\sigma_{10}(p)\in C\}
\qquad
\bigl(C\subset S_{10}\ \text{为共轭类}\bigr)
\]
生成的有限布尔代数，令 \(\mathcal A_3\) 为由全部事件
\[
\{\sigma_3(p)\in D\}
\qquad
\bigl(D\subset S_3\ \text{为共轭类}\bigr)
\]
生成的有限布尔代数，并记
\[
\chi_{10}(p):=\mathrm{sgn}\!\bigl(\sigma_{10}(p)\bigr).
\]
则对任意
\[
E\in\mathcal A_{10},
\qquad
F\in\mathcal A_3,
\]
都有精确乘法密度律
\[
\delta(E\cap F)=\delta(E)\,\delta(F).
\]
特别地，凡是仅由 \(A\)、\(\{\chi_{10}=+1\}\)、\(\{\chi_{10}=-1\}\) 及其他 \(S_{10}\)-循环型事件经有限次并、交、补得到的事件，均与仅由 \(B_{\mathrm{split}}\)、\(B_{\mathrm{single}}\)、\(B_{\mathrm{noroot}}\)、\(B_{\mathrm{root}}\)、\(\{\chi_{\mathrm{LY}}=+1\}\)、\(\{\chi_{\mathrm{LY}}=-1\}\) 及其他 \(S_3\)-分裂型事件经有限次并、交、补得到的事件严格独立。
\end{theorem}
\begin{proof}
由有限群上的布尔代数结构，每个 \(E\in\mathcal A_{10}\) 都可唯一写成若干 \(S_{10}\) 共轭类的有限不交并：
\[
E=\{\sigma_{10}(p)\in \widetilde E\},
\qquad
\widetilde E\subset S_{10}\ \text{为共轭类并},
\]
同理每个 \(F\in\mathcal A_3\) 都可写成
\[
F=\{\sigma_3(p)\in \widetilde F\},
\qquad
\widetilde F\subset S_3\ \text{为共轭类并}.
\]
因此群上的指示函数
\[
f:=\mathbf 1_{\widetilde E},
\qquad
h:=\mathbf 1_{\widetilde F}
\]
都是类函数。把它们代入定理
\ref{thm:xi-time-part62ea-dualfingerprint-classfunction-prime-average-factorization}，得到
\[
\delta(E\cap F)
=
\left(\frac{|\widetilde E|}{10!}\right)
\left(\frac{|\widetilde F|}{6}\right)
=
\delta(E)\,\delta(F).
\]
这就证明了强乘法律。

最后，\(A\) 是 \(S_{10}\) 中“存在不动点”所对应的共轭类并；\(\{\chi_{10}=\pm1\}\) 是偶/奇置换所对应的共轭类并；而 \(B_{\mathrm{split}},B_{\mathrm{single}},B_{\mathrm{noroot}},B_{\mathrm{root}},\{\chi_{\mathrm{LY}}=\pm1\}\) 全部都是 \(S_3\) 共轭类事件或其并。故列举出的有限布尔组合都落在相应布尔代数中，从而自动满足上式。证毕。
\end{proof}

\endinput
